\chapter{Fonctions linéaires et affines}
\label{exercices:09}
\section*{Exercice 9.1 -- Identifier et tracer}
Parmi $f(x)=3x$, $g(x)=-2x+4$, $h(x)=5$ et $u(x)=x^2$, lesquelles sont linéaires, affines, constantes ? Tracer $f$ et $g$ sur un même repère.

\section*{Exercice 9.2 -- Retrouver une expression}
Une fonction affine $f$ vérifie $f(1)=5$ et $f(4)=14$. Déterminer $f(x)$, puis l'antécédent de $-1$.

\section*{Exercice 9.3 -- Interpréter les paramètres}
Une température modélisée en fonction du temps vaut $T(t)=22-0{,}4t$, en degrés Celsius, pour $0\leq t\leq30$ minutes. Interpréter $22$ et $-0{,}4$, calculer $T(15)$ et le temps pour atteindre $12\degre$C.

\section*{Exercice 9.4 -- Comparer deux abonnements}
Les coûts pour $n$ utilisations sont $A(n)=20+2n$ et $B(n)=5+3n$. Chercher l'égalité des prix, puis l'offre la moins chère selon le nombre entier $n\geq0$.
